% Supplementary Material for manuscript AESM-2026-0104 (revised).
% All tables are generated from the stored result files by
% Code/revision_analysis/rev_supplement_tables.py; no value is typed by hand.
\documentclass[10pt,a4paper]{article}
\usepackage[margin=2.2cm]{geometry}
\usepackage[T1]{fontenc}
\usepackage[utf8]{inputenc}
\usepackage{booktabs,array,graphicx,amsmath}
\usepackage[hidelinks]{hyperref}
\renewcommand{\thetable}{S\arabic{table}}
\renewcommand{\thesection}{S\arabic{section}}
\newcommand{\T}[1]{\input{supplement_tables/#1}}
\setlength{\parskip}{4pt}\setlength{\parindent}{0pt}
\title{Supplementary Material\\[4pt]\large Quantum Machine Learning for Cyber-Physical Anomaly Detection in
Unmanned Aerial Vehicles: Lessons from a Proxy-Audited Temporal Evaluation}
\author{C.~A.~Dur\'an~Paredes, J.~E.~Le\'on~Calder\'on, N.~S\'anchez~Perea, G.~D.~D\'iaz, and C.~Segura}
\date{}
\begin{document}
\maketitle

This document holds the detailed tables behind the article. Unless noted,
values are means with standard deviations in parentheses over the ten
partitions (``seeds'') of the same log. The standard deviation measures
sensitivity to the choice of partition. It is not a confidence interval for
performance on an independent flight. \emph{all}, \emph{reduced}, and
\emph{physical} denote the three input sets (72, 63, and 22 variables).
The false-alarm rate is the fraction of nominal samples flagged as faults;
the miss rate is the fraction of fault samples labeled nominal. All models
use default decision thresholds.

\section{Data, labels, and partitions}
The table rebuilt from the raw sensor files has 4\,817 rows (about 10\,Hz)
and 81 numeric variables, of which nine are constant and were dropped. The
timestamps run on one clock for 506\,s. Two logging gaps of 16.3\,s and
8.4\,s separate three segments of 144, 4\,606, and 67 rows; the first
segment is entirely nominal and the last entirely remote-control fault.
There are four contiguous fault intervals: label~1 from 92.0 to 152.7\,s
(608 rows), label~2 from 226.0 to 285.3\,s (594 rows), label~3 from 335.0
to 383.1\,s (482 rows), and label~4 from 430.9 to 506.1\,s (670 rows).

\emph{Check of the label meanings.} The data set's authors describe four
simulated failures in the order GPS, accelerometer, engine, and remote
control. The raw signals agree with labels 1 to 4 in that order. During
label~1 the reported satellite count drops from 10 to 3. During label~2
all three accelerometer channels are constant. During label~3 the ground
speed is zero, the altitude estimate is at ground level, and the throttle
output is saturated. During label~4 the vehicle descends.

Table~\ref{tab:blocks} gives the ten time blocks and Table~\ref{tab:parts}
the ten partitions. In every partition some training rows are later in
time than the earliest test row, and at least one test block borders a
training block. Blocks 0 and 7 are never used for testing; block 8 is used
in six partitions. In seven partitions at least one fault interval in the
test blocks contributes no rows to training. The median absolute
autocorrelation of the 22 physical variables is 0.87 at 0.1\,s, 0.67 at
1\,s, 0.46 at 10\,s, and 0.07 at 48\,s (one block).

\begin{table}[h]\centering\small
\caption{The ten time blocks and their label counts.}\label{tab:blocks}
\T{S2_blocks}
\end{table}

\begin{table}[h]\centering\footnotesize
\caption{The ten partitions. Validation blocks are the two blocks not listed.}\label{tab:parts}
\resizebox{\linewidth}{!}{\T{S2_partitions}}
\end{table}

\section{Complete results of the blocked evaluation}
Tables~\ref{tab:f1}--\ref{tab:miss} give all twelve models and all three
input sets. In the stored result files the column \texttt{far\_class0} is
the miss rate; the false-alarm rate was recovered from the same files as
$2\times$\texttt{far\_macro}$-$\texttt{far\_class0} and checked against
recomputed confusion matrices for the classical models (maximum difference
$2\times10^{-16}$).

\begin{table}[h]\centering\small\caption{F1 score (macro average).}\label{tab:f1}\T{S1_f1}\end{table}
\begin{table}[h]\centering\small\caption{Area under the ROC curve.}\label{tab:auc}\T{S1_auc}\end{table}
\begin{table}[h]\centering\small\caption{False-alarm rate.}\label{tab:far}\T{S1_far}\end{table}
\begin{table}[h]\centering\small\caption{Miss rate.}\label{tab:miss}\T{S1_miss}\end{table}

Table~\ref{tab:wins} compares the trained-circuit hybrid with every other
model partition by partition (physical inputs). Each entry counts the
partitions in which the hybrid is better / tied / worse.

\begin{table}[h]\centering\small
\caption{Trained-circuit hybrid against each other model: partitions won / tied / lost (physical inputs).}\label{tab:wins}
\T{S5_wins}
\end{table}

\clearpage
\section{Effect of the split design}
Table~\ref{tab:split} repeats the classical experiments with a random
stratified split of rows (70/15/15, ten seeds) and with the blocked split
after deleting every training row within a given time gap of any test row.
The average number of real training rows falls from 2\,890 (no gap) to
2\,765, 2\,515, 2\,152, and 1\,734, so the gap analysis confounds reduced
adjacency with reduced training data.

\begin{table}[h]\centering\small
\caption{Mean F1 of the classical models under different splits.}\label{tab:split}
\T{S3_split_design}
\end{table}

\emph{Training on the past only.} Models were trained on blocks $0,\dots,k-1$
and tested on block $k$ for $k=2,\dots,9$. Table~\ref{tab:fwd} gives the
accuracy per test block with physical inputs and Table~\ref{tab:fwdpool}
the pooled error rates. On block~4, the first block of the accelerometer
fault, no model detects a single fault sample.

\begin{table}[h]\centering\small
\caption{Past-only training: accuracy on each test block (physical inputs).}\label{tab:fwd}
\T{S3_forward}
\end{table}
\begin{table}[h]\centering\small
\caption{Past-only training: results pooled over test blocks 2--9.}\label{tab:fwdpool}
\T{S3_forward_pooled}
\end{table}

\emph{Holding out gap-delimited segments} is not informative: the two short
segments contain a single class, and the long one leaves 211 training rows.

\section{Detection by exposure to the fault event}
For each partition, test fault samples were divided according to whether
their fault interval also contributed rows to training. Table~\ref{tab:events}
pools the ten partitions (14\,745 and 11\,645 sample-level decisions per
input set across the five models).

\begin{table}[h]\centering\small
\caption{Fraction of fault samples detected.}\label{tab:events}
\T{S4_events}
\end{table}

\clearpage
\section{What each model received, and controls}
\emph{Classical models:} all variables of the input set; all training rows
after rebalancing (about 2\,900--3\,800, of which 2\,890 are real); default
settings; fitted in under two seconds. \emph{Circuit:} the five variables
ranked highest by mutual information on the training blocks; subset~A, 800
rows (400 per class); 30 weights; 200 optimizer iterations.
\emph{Hybrids and controls:} the same five variables plus 0 to 20 derived
columns; subset~B, 800 rows disjoint from A; XGBoost with 100 trees.
\emph{Trained-circuit hybrid:} its extra column comes from a circuit fitted
on the labels of subset~A, so it draws on 1\,600 labeled rows, against 800
for the controls. The validation blocks were not used. Subsets A and B are
drawn from the rebalanced training set, in which 12\% of rows on average
are synthetic, so they are disjoint but not independent.

``Identical budget'' in the original submission meant the same partitions
and test blocks for all models, and the same subset~B and XGBoost settings
for the six hybrids. It did not mean equal samples, inputs, computation, or
capacity across model families.

Table~\ref{tab:fair} traces the difference between classical XGBoost and
the ``XGB only'' control step by step, and adds two label-informed
classical controls. With all variables, the five selected variables are
accumulating quantities whose test values fall outside the training range
in 30\% of test rows on average; the ``XGB only'' control then predicts a
single class in seven of ten partitions.

\begin{table}[h]\centering\small
\caption{Mean F1 (standard deviation) for the fairness decomposition and added controls.}\label{tab:fair}
\T{S5_fairness}
\end{table}

\section{Sensitivity to synthetic oversampling}
SMOTETomek was applied after the split, to training rows only.
Table~\ref{tab:res} repeats the classical experiments without it.
Table~\ref{tab:d2} repeats the circuit and hybrid experiments on real rows
only (physical inputs; ten new circuit fits). Table~\ref{tab:syn}
characterizes the synthetic rows. On average 9\% of them (up to 23\%) have
a magnetic-field magnitude outside the central 98\% range of the real
training rows, where 2\% would be expected. We make no claim that the
synthetic samples are physically plausible.

\begin{table}[h]\centering\footnotesize
\caption{Classical models with and without SMOTETomek (means over ten partitions).}\label{tab:res}
\T{S6_resampling}
\end{table}
\begin{table}[h]\centering\footnotesize
\caption{Circuit and hybrids without synthetic rows (physical inputs).}\label{tab:d2}
\resizebox{\linewidth}{!}{\T{S6_dru_no_resampling}}
\end{table}
\begin{table}[h]\centering\footnotesize
\caption{Synthetic rows in the rebalanced training sets (physical inputs).}\label{tab:syn}
\resizebox{\linewidth}{!}{\T{S6_synthetic}}
\end{table}

\clearpage
\section{Cost of the simulated circuit}
Table~\ref{tab:cost} reports measurements on the classical computer that
ran the simulation. They do not describe a quantum processor. On hardware,
each expectation value would be estimated from repeated circuit runs
affected by gate and readout noise, with queueing and communication delays;
none of these was measured. The 30 circuit fits of the original study took
117 minutes in total. A second measurement of the same fit gave 234\,s and
1.4\,ms per sample; all other values were identical.

\begin{table}[h]\centering\small
\caption{Measured cost of one circuit fit and of inference in simulation.}\label{tab:cost}
\T{S7_cost}
\end{table}

\section{Secondary experiment on the engine fault}
A one-fault task (nominal versus engine fault) was run with twenty time
blocks and XGBoost. In nine of the ten random assignments at least one of
the training, validation, and test sets contained no fault sample, and
those assignments were skipped; one remained (seed~1). On that single
partition F1 was 0.93 with the \emph{all} and \emph{reduced} inputs and 0.43
with \emph{physical} inputs, where every test row was labeled nominal
although the ROC AUC stayed at 0.96. A single partition supports no
estimate and is reported only as an anecdote.

\section{Reproducibility}
Every stored result that was rerun was reproduced exactly: all 150
classical runs, the 30 runs of the five-variable XGBoost control, and one
circuit fit (seed~0, \emph{physical} inputs; Table~\ref{tab:cost}). The other
29 circuit fits of the main study, and the hybrid variants that depend on
them, were not rerun. The circuit and hybrid rows of
Tables~\ref{tab:f1}--\ref{tab:miss} and the comparisons in
Table~\ref{tab:wins} are therefore computed from the stored per-partition
files of the original study; the false-alarm rate in those rows is derived
arithmetically from two stored columns. All other tables come from analyses
executed for this revision. Scripts, per-partition outputs, software versions,
and commands are in the public repository
\url{https://github.com/Carlosandp/TLM-UAV-Quantum-Anomaly-Detection}
(folder \texttt{revision/}).
\end{document}
